% Part: first-order-logic
% Chapter: natural-deduction
% Section: rules-and-proofs

\documentclass[../../../include/open-logic-section]{subfiles}

\begin{document}

\iftag{FOL}
      {\olfileid[pt]{fol}{ntd}{rul}}
      {\olfileid[pt]{pl}{ntd}{rul}}

\olsection{Regras e \usetoken{P}{derivation}}

\begin{explain}
Sistemas de dedução natural pretendem ser um paralelo próximo do raciocínio
informal usado em provas matemáticas (daí serem, de certo modo,
``naturais''). Provas em dedução natural começam com hipóteses.
Em seguida, aplicam-se regras de inferência. As hipóteses são ``!!{discharged}s''
pelas regras de inferência \Intro{\lnot}, \Intro{\lif}, \iftag{FOL}{\Elim{\lor} e
  \Elim{\lexists}}{ e \Elim{\lor}}, e o rótulo da
hipótese !!{discharged} é colocado ao lado da inferência para maior
clareza.
\end{explain}

\begin{defn}[Hipótese]
Uma \emph{hipótese} é qualquer !!{sentence}
na posição mais alta de qualquer ramo.
\end{defn}

!!^{derivation}s em dedução natural são certas árvores de
!!{sentence}s, em que as !!{sentence}s no topo são hipóteses e, se
!!a{sentence} está abaixo de um, dois ou três outros sequentes, ela deve
seguir corretamente por uma regra de inferência. As !!{sentence}s no topo
da inferência são chamadas de \emph{premissas} e a !!{sentence}
abaixo, de \emph{conclusão} da inferência.  As regras vêm em
pares, uma regra de introdução e uma de eliminação para cada
!!{operator}. Elas introduzem !!a{operator} na conclusão ou
removem !!a{operator} de uma premissa da regra.  Algumas das regras
permitem que uma hipótese de certo tipo seja \emph{!!{discharged}}. Para
indicar qual hipótese é !!{discharged} por qual inferência, também
atribuímos rótulos tanto à hipótese quanto à inferência.  Isso é
indicado escrevendo-se a hipótese como ``$\Discharge{!A}{n}$.''

% Only include this sentence is on of \land, lor, lif or lnot is defined.

\iftag{notprvNot,notprvAnd,notprvOr,notprvIf}{}%
	{É costume considerar regras para todos os !!{operator}s $\land$, $\lor$, $\lif$, $\lnot$ e $\lfalse$, mesmo que alguns deles sejam definidos.}



\end{document}
